\section*{摘要}\label{ux6458ux8981}

早期眼底筛查是眼科疾病诊断的有效方法。近年来，深度学习的发展为眼底图像分析提供了强大工具。然而，真实临床部署仍面临三个层层递进的挑战：（1）单只眼睛中常同时存在多种疾病，若将各疾病独立建模，容易造成误检与漏检；（2）跨机构部署存在数据分布偏移，而常见的跨域方法通常预设可访问目标域或多源域数据，这一前提在真实医疗场景中难以满足；（3）医疗数据具有较强的隐私敏感性，进一步限制了跨机构数据共享，使得多源域的预设难以成立，因此需要在本地隐私保护前提下实现良好的泛化性能。针对上述三个挑战，我们首先提出自适应标签-特征关联性构建模块（Adaptive label-feature Relevance Construction, ARC），通过 Transformer 注意力机制捕捉疾病共现依赖，减少多标签识别中的误检与漏检；其次设计双分支引导域不变特征提取模块（Dual-branch guided domain-invariant Feature Extraction, DFE），在类激活映射（Class Activation Mapping, CAM）引导下仅使用单一源域数据学习域不变特征，在提升未知域泛化性能的同时增强模型可解释性；最后将自适应高斯差分隐私（Adaptive Gaussian Differential Privacy, AdaGDP）机制集成至训练流程中，通过动态梯度噪声注入实现隐私保护，并将性能损失控制在可接受范围。实验表明，所提方法在 ODIR 上取得 78.85 的 Final score，优于所选代表性多标签眼科疾病识别方法，并在 EDID 跨域数据集上取得 3.19 个百分点的提升。综上，本文搭建了一个面向多标签眼科疾病识别的可解释与隐私保护单域泛化框架，提升了跨域泛化性能，并为隐私保护场景下的眼科疾病诊断提供启发。

\noindent\textbf{关键词：}\textbf{Domain Generalization; Multi-Label Ocular Disease Recognition; Data Privacy; Interpretability}
